% Prisma — Early Results Paper (F1-AC-001 / F1-AC-002).
%
% Status: DRAFT skeleton, weeks-1-of-Fase-1. Outline + section
% scaffolding only. Numbers, figures, and the related-work section
% land as the implementation matures.

\documentclass[11pt,a4paper]{article}

\usepackage[margin=1in]{geometry}
\usepackage{microtype}
\usepackage{hyperref}
\usepackage{listings}
\usepackage{graphicx}
\usepackage{booktabs}

\title{Prisma: A Provably-Sound, NPU-Assisted DBT for x86_64 on ARM64}
\author{Danny Gonz\'{a}lez Ram\'{i}rez \\ \texttt{danny@prisma-emu.dev}}
\date{Draft --- \today}

\begin{document}
\maketitle

\begin{abstract}
We present \emph{Prisma}, a dynamic binary translator for the
x86\_64 \(\rightarrow\) ARM64 path that targets three goals
simultaneously: \textbf{(1)} mechanically-verified pass soundness
via Lean 4, \textbf{(2)} adaptive memory-model
\(\textsc{tso}\)\(\rightarrow\)relaxed rewriting backed by the
proof, and \textbf{(3)} an on-device NPU-assisted optimisation
classifier that decides per-region between fast interpretation and
deeper compilation. Early results from a clean-slate IR + lowering
pipeline targeting Apple silicon and Adreno-class Android show
\emph{[NUMBERS]} on Dhrystone and \emph{[NUMBERS]} on coreutils
relative to QEMU user-mode. We outline the proof obligations
already discharged in Lean and the path to closing the remaining
sorries.
\end{abstract}

\section{Introduction}
\label{sec:intro}

Dynamic binary translation has carried thirty years of value
movement across ISAs: Transmeta Crusoe, Intel IA-32 EL, Apple
Rosetta 2, FEX-Emu and Box64. Each chose a sweet spot on the
correctness-vs-performance frontier and lived with the tail of
guest binaries that exposed the cracks. Prisma takes a deliberately
maximalist position: \emph{provable} pass soundness, \emph{adaptive}
memory-model rewriting, and \emph{ML-assisted} per-region
compilation strategy. The cost is up-front engineering (a Lean 4
spec of the IR before any pass ships); the payoff is a system that
can be cited and rebuilt.

\paragraph{Contributions.}
\begin{itemize}
  \item A clean-slate SSA IR with 32 opcodes, formally specified
        in Lean 4 (\texttt{ir-spec/}).
  \item A complete pass pipeline (constant prop, algebraic
        simplification, CSE, copy prop, redundant-load /
        dead-store elimination, branch fold, peephole, tail-call,
        DCE) with one published soundness proof and three open
        proof obligations (\S\ref{sec:proofs}).
  \item An ARM64 lowering that integrates W\^{}X-safe MAP\_JIT
        slabs, page-protection-based SMC detection (RFC 0010),
        and a separated FP / SIMD register pool with an explicit
        Flags pillar (RFC 0004).
  \item An end-to-end measurement of Dhrystone, CoreMark, and
        a coreutils subset against QEMU user-mode and Box64.
\end{itemize}

\section{IR Design}
\label{sec:ir}

\subsection{SSA, sized, and typed}
Refs are \texttt{uint32\_t}. Every \texttt{Ref} carries an
\texttt{OpSize} (I8/I16/I32/I64) tracked by the validator
(F1-IR-015). The validator catches BinOp / Compare / StoreReg size
mismatches before they reach the lowerer.

\subsection{Opcodes}
The 32-opcode surface (RFC 0001, RFC 0009) covers ALU, memory
(plain + TSO), control flow (block-id and guest-PC variants),
explicit Flags ops (RFC 0004), scalar floating-point, 128-bit SIMD
support in the emitter, and an \texttt{InlineAsm} escape hatch
for the long tail of x86 instructions not yet decoded.

\subsection{Serialisation}
The IR has a stable on-wire format (RFC 0009) so the persistent
translation cache can stash post-passes IR for re-lowering when a
host CPU feature flag changes.

\section{Lowering and Backend}
\label{sec:backend}

The lowerer maps SSA Refs to host registers via two pools: \\
\(\bullet\) integer pool (x0..x9 scratch + x10..x26 pinned guest
GPRs); \\
\(\bullet\) FP pool (V0..V7 scratch). \\
The Flags pillar uses a separate ``flag\_refs\_'' set instead of a
register, since flags live in NZCV. Dominator-based natural-loop
analysis (F1-IR-024 / F1-IR-025) feeds the planned register
allocator (RFC 0006).

\section{Memory Model}
\label{sec:memory}

x86 enforces TSO; ARM relaxes. Prisma emits \texttt{ldar} /
\texttt{stlr} for every guest load / store by default
(\texttt{LoadMemTSO} / \texttt{StoreMemTSO}). The TSO-adaptive
pass (Pillar 3, Fase 2.5) rewrites these to plain \texttt{ldr} /
\texttt{str} when proven safe; the proof obligation is the
RFC's \emph{TSO-rewrite soundness} statement.

\section{Proofs}
\label{sec:proofs}

The Lean 4 spec defines \texttt{Op}, \texttt{Stmt}, \texttt{BasicBlock},
\texttt{Function}, \texttt{evalBinOp}, and a small
\texttt{MachineState} with the GPR file. Three lemmas are
discharged today: \texttt{evalBinOp\_determinism},
\texttt{constant\_reduction}, \texttt{binop\_reduction}. One
\texttt{sorry} remains in \texttt{Lemmas.lean}, gated by CI.

\section{Implementation}
\label{sec:impl}

\subsection{Languages and tools}
C++20 (core engine, \texttt{vixl} ARM64 assembler), Rust 2024 (PE
loader and shell, planned), Kotlin 2.0 + Jetpack Compose (Android
app, planned), Lean 4 (spec + proofs), Python 3.12 (build / bench
tooling).

\subsection{Test surface}
537 Catch2 test cases / 3783 assertions at the time of writing,
covering decoder, IR, passes, lowering, runtime, and cache.
Property-based tests (rapidcheck) and a differential harness
against QEMU user-mode are planned (F1-TC-003 / F1-TC-010).

\section{Early Measurements}
\label{sec:eval}

\emph{[Placeholder until F1-AC-003 / F1-AC-004 land. Target table:
Dhrystone MIPS / CoreMark / coreutils geometric-mean speedup vs
QEMU user-mode and Box64 on Apple M3 Max, Snapdragon 8 Gen 3,
Dimensity 8300-Ultra.]}

\section{Related Work}
\label{sec:related}

\emph{[Placeholder. Tier 1: FEX-Emu, Box64. Tier 2: Apple Rosetta 2,
Intel IA-32 EL, Transmeta Crusoe. Tier 3: QEMU TCG, ELITE,
DynamoRIO. CompCert / CakeML for the proof tradition. TOSTING for
TSO-rewrite quantitative work.]}

\section{Conclusion and Future Work}
\label{sec:conc}

Prisma is a long-term bet on three pillars rarely combined:
provable correctness, adaptive memory-model relaxation, and
per-region NPU-assisted strategy choice. This early-results draft
documents the IR, lowering, and proof scaffolding that exist
today; the deep evaluation (\S\ref{sec:eval}) and TSO-adaptive
proof are the next milestones.

\bibliographystyle{plain}
\bibliography{refs}

\end{document}
